\documentclass[12pt]{article}
\usepackage{sbc-template}
\usepackage[utf8]{inputenc}
\usepackage[T1]{fontenc}
\usepackage[brazilian]{babel}
\usepackage{graphicx,url}
\usepackage{listings}
\sloppy

% Data do documento solicitada pelo autor; não representa a data dos testes.
\newcommand{\datadotrabalho}{20 de outubro de 2026}
\date{\datadotrabalho}
\lstset{language=C,basicstyle=\ttfamily\footnotesize,
  columns=fullflexible,keepspaces=true,showstringspaces=false,
  breaklines=true,tabsize=4,aboveskip=6pt,belowskip=6pt}
\title{Sistema de Arquivos com FUSE: Desenvolvimento Parcial de um Arquivo Virtual de Monitoramento}
\author{Mateus Borges Rosenzweig\inst{1}}
\address{Engenharia de Computação -- Universidade Federal do Ceará\\
Quixadá -- Ceará -- Brasil
\email{mateus.rosenzweig@alu.ufc.br}}

\begin{document}
\maketitle
\begin{center}
\small \datadotrabalho
\end{center}

\begin{abstract}
This paper reports the partial development of a FUSE filesystem based
on the libfuse passthrough example. A read-only virtual file, status.txt,
reports directory creations, file openings, and read requests. Its contents
are generated in memory for each opening. Direct callback tests checked
partial reads, concurrent counters, directory listing, and write protection.
The final stage will add byte counters and a virtual history of recent
operations, followed by integration tests through a real FUSE mount.
These extensions have not yet been implemented.
\end{abstract}

\begin{resumo}
Este trabalho apresenta o desenvolvimento parcial de um sistema de
arquivos com FUSE baseado no exemplo passthrough da libfuse. O arquivo
virtual status.txt informa criações de diretórios, aberturas e pedidos de
leitura, com conteúdo gerado em memória a cada abertura. Testes diretos
dos callbacks verificaram leitura parcial, concorrência, listagem e
proteção contra escrita. Para a entrega final, são propostas a contagem
de bytes e a consulta a um histórico virtual de operações, além de testes
por uma montagem real. Essas extensões ainda não foram implementadas.
\end{resumo}

\section{Introdução}

Sistemas de arquivos permitem que aplicações acessem dados por meio de
nomes, diretórios e operações conhecidas, como abrir, ler e escrever.
Essas operações fazem parte do uso cotidiano do computador, mas os
programas que as utilizam geralmente não precisam conhecer a forma como
os dados são armazenados. Estudar essa interface ajuda a compreender a
relação entre aplicações e sistema operacional.

A atividade foi proposta pela disciplina de Projeto e Implementação
de Sistemas Operacionais, com o objetivo de entender e implementar um
sistema de arquivos com funcionalidades próprias. Para isso, foi utilizada
a interface FUSE, que permite atender a operações de arquivos em um
programa executado em espaço de usuário. A biblioteca libfuse auxilia
na comunicação desse programa com o componente FUSE do kernel
\cite{libfuse}.

O objetivo geral do projeto é adaptar o exemplo \texttt{passthrough.c}
para disponibilizar informações sobre seu funcionamento por meio de
arquivos virtuais. Esta entrega parcial apresenta a primeira etapa,
cuja funcionalidade implementada é o arquivo
\texttt{status.txt}, acessível na raiz do ponto de montagem. Ele apresenta
três contadores: diretórios criados com sucesso, aberturas de arquivos
reais bem-sucedidas e pedidos de leitura recebidos pelo callback.

A funcionalidade pode ser observada com comandos comuns, como
\texttt{cat}. Nesta etapa, o foco é compreender os callbacks e sua
integração. A entrega final deverá ampliar o monitoramento com volume
de dados transferidos e histórico de operações, conforme o plano da
Seção~\ref{sec:plano-final}.

\section{Fundamentação Teórica}

\subsection{Arquivos, diretórios e operações}

Um arquivo possui conteúdo e metadados, como tipo, tamanho e permissões.
Diretórios organizam entradas por nomes. Consultar atributos, listar
entradas e ler conteúdo são operações distintas nessa interface.

Quando um arquivo real é aberto pelo processo, a função \texttt{open}
retorna um descritor. Esse número identifica a abertura e pode ser usado
em operações posteriores. No passthrough, o descritor é guardado no campo
\texttt{fh} de \texttt{fuse\_file\_info}. A implementação pode, assim,
reutilizá-lo ao atender leituras e fechar o recurso ao final
\cite{passthrough}.

\subsection{FUSE e a biblioteca libfuse}

FUSE significa \textit{Filesystem in Userspace}. Seu funcionamento
envolve um componente no kernel e um programa em espaço de usuário.
A libfuse oferece uma API de alto nível, baseada em caminhos, e uma API
de baixo nível, baseada em inodes. Este projeto utiliza a primeira
\cite{libfuse}.

O ponto de montagem é o diretório pelo qual o sistema de arquivos fica
acessível. Quando uma aplicação acessa um caminho dentro dele, o kernel
pode encaminhar pedidos ao programa FUSE. A libfuse chama as funções
registradas em \texttt{fuse\_operations}, chamadas de callbacks. O programa
define como cada pedido será atendido e devolve o resultado.

Um callback pode ter um nome escolhido pelo programador, mas deve
respeitar a assinatura da operação. Por exemplo, o campo \texttt{.mkdir}
pode apontar para \texttt{xmp\_mkdir} ou para outra função compatível.
Isso permite alterar o comportamento de criação de diretórios sem
inventar uma nova operação no protocolo. Na API de alto nível, erros
são normalmente devolvidos como valores negativos, como
\texttt{-ENOENT} para um caminho inexistente \cite{operacoes}.

\subsection{O exemplo passthrough}

O exemplo \texttt{passthrough.c} repassa solicitações para funções que
operam no sistema de arquivos existente. Entre elas estão
\texttt{lstat}, \texttt{mkdir}, \texttt{pread} e \texttt{pwrite}.
Na forma utilizada, ele espelha a hierarquia iniciada em \texttt{/}
\cite{passthrough}. Assim, se estiver montado em \texttt{/tmp/meufs},
o caminho \texttt{/tmp/meufs/etc/hostname} corresponde ao arquivo real
\texttt{/etc/hostname}.

Essa base permite concentrar o estudo nos callbacks. Alterações nos
caminhos espelhados podem modificar arquivos reais, conforme as permissões
do processo. O monitoramento foi adicionado como exceção ao encaminhamento.

\section{Implementação da Etapa Parcial}

\subsection{Organização do arquivo virtual}

O caminho reservado é \texttt{/status.txt}, considerando a raiz do
sistema montado. A função auxiliar \texttt{is\_status} compara o caminho
recebido com esse nome. Quando há correspondência, o callback atende
ao arquivo virtual. Nos demais casos, continua o fluxo do passthrough.

Esse arquivo não é criado no disco. Seu texto é produzido pela função
\texttt{make\_status}, que usa \texttt{snprintf} para reunir os contadores.
A estrutura \texttt{status\_snapshot} guarda o texto em um vetor de
256 bytes e registra seu comprimento. Esse espaço atende às três
mensagens e aos valores dos contadores utilizados.

A Tabela~\ref{tab:callbacks} resume as principais adaptações. Apenas
implementar a leitura seria insuficiente: o caminho também precisa
ser reconhecido na consulta de atributos e aparecer na listagem.

\begin{table}[ht]
\centering
\caption{Participação dos callbacks no arquivo virtual.}
\label{tab:callbacks}
\begin{tabular}{p{0.20\textwidth}p{0.71\textwidth}}
\hline
Callback & Adaptação realizada \\
\hline
\texttt{getattr} & Informa arquivo regular, tamanho e permissões 0444. \\
\texttt{readdir} & Inclui o nome na raiz da montagem. \\
\texttt{open} & Aceita leitura e guarda uma cópia do texto. \\
\texttt{read} & Copia o intervalo solicitado para o buffer. \\
\texttt{release} & Libera a memória associada à abertura. \\
\texttt{write} & Recusa alterações no conteúdo virtual. \\
\hline
\end{tabular}
\end{table}

Na listagem da raiz, o nome é inserido por meio de \texttt{filler}.
Se houver um arquivo real com o mesmo nome nessa raiz, sua entrada
é omitida para evitar duplicação. Essa escolha reserva o nome para o
monitoramento. Um arquivo chamado \texttt{status.txt} dentro de outro
diretório continua seguindo o comportamento normal.

\subsection{Contagem das operações}

Os contadores foram declarados com o tipo
\texttt{atomic\_uint\_fast64\_t}. O uso de operações atômicas evita a
perda de incrementos quando mais de uma thread atende solicitações.
O contador de diretórios é incrementado somente depois de
\texttt{mkdir} ter sucesso. O trecho abaixo foi extraído dessa parte
do callback:

\begin{lstlisting}
res = mkdir(path, mode);

if (res == -1)
    return -errno;

atomic_fetch_add(&directories_created, 1);
return 0;
\end{lstlisting}

Aberturas reais são contadas após o sucesso nos callbacks
\texttt{open} e \texttt{create}. Já o contador de leitura é incrementado
na entrada do fluxo de arquivos reais em \texttt{xmp\_read}, incluindo
pedidos que terminem em erro. Essa diferença precisa ser considerada
na interpretação dos números: dois contadores representam sucessos
e o terceiro representa solicitações ao callback.

O próprio arquivo virtual é tratado antes dos incrementos de abertura
e leitura. Consultar \texttt{status.txt}, portanto, não aumenta esses
contadores. Os dados permanecem apenas em memória e são reiniciados
quando o processo do sistema de arquivos é iniciado novamente.

\subsection{Abertura, leitura e liberação}

A cada abertura do arquivo virtual, o programa aloca uma
\texttt{status\_snapshot} e preenche seu texto. O endereço da estrutura
é armazenado em \texttt{fi->fh}. Nesse caso, esse campo guarda um
ponteiro convertido para o tipo inteiro aceito pela API, enquanto
para os arquivos reais continua guardando um descritor.

A decisão de manter uma cópia por abertura evita misturar trechos de
textos produzidos em momentos diferentes. Se outra thread atualizar
os contadores durante uma leitura, a abertura atual continuará
utilizando sua cópia. Para observar novos valores, o leitor deve
fechar e abrir o arquivo novamente.

A leitura considera \texttt{size}, que é a quantidade solicitada,
e \texttt{offset}, que indica a posição inicial. Depois de obter a
cópia correspondente à abertura, o callback executa:

\begin{lstlisting}
if (offset < 0)
    return -EINVAL;
if ((uintmax_t)offset >= snapshot->length)
    return 0;
if (size > snapshot->length - (size_t)offset)
    size = snapshot->length - (size_t)offset;
memcpy(buf, snapshot->text + (size_t)offset, size);
return (int)size;
\end{lstlisting}

O retorno zero informa o fim do arquivo. Limitar o tamanho da cópia
evita acessar posições além do texto. Cada pedido recebe apenas o
intervalo disponível a partir da posição solicitada.

A abertura virtual também configura \texttt{direct\_io = 1} e
\texttt{keep\_cache = 0}. Isso permite que novas aberturas consultem
o texto atualizado, sem reutilizar o cache de dados do arquivo.
Ao encerrar a abertura, \texttt{release} libera a memória com
\texttt{free}. Usar \texttt{close} nesse caso seria incorreto, pois
o recurso guardado é uma estrutura em memória.

\subsection{Proteção e tratamento de erros}

Os atributos apresentam o arquivo como somente leitura, mas a
implementação também recusa operações que poderiam alterá-lo.
Escrita, truncamento, remoção e renomeação envolvendo o caminho
reservado são bloqueados nos callbacks. A criação de outra entrada
nesse mesmo caminho retorna um erro de existência.

Também foram tratados os callbacks opcionais presentes na base,
incluindo atributos estendidos, consulta com \texttt{statx} e cópia
por intervalo. Quando o arquivo virtual é a origem de
\texttt{copy\_file\_range}, a função retorna operação não suportada,
permitindo que um cliente compatível tente copiar por leitura e
escrita. O objetivo é evitar que essas rotas usem o recurso virtual
como se ele fosse um descritor de arquivo real.

\section{Metodologia da Avaliação Parcial}

A implementação foi compilada com GCC 16.2.1 e libfuse 3.18.3.
O código de referência pertence ao repositório libfuse, na revisão
\texttt{fe66006f}, com as modificações locais descritas neste artigo.
Foram verificadas compilações com e sem as opções de callbacks
condicionais, incluindo \texttt{HAVE\_STATX} e
\texttt{HAVE\_COPY\_FILE\_RANGE}.

A avaliação foi feita por um programa de testes em C que inclui o
arquivo da implementação e chama seus callbacks diretamente.
As operações reais utilizaram um diretório temporário. Para verificar
a listagem da raiz e a colisão com um nome existente, o teste
redirecionou a chamada de abertura do diretório raiz para uma pasta
de teste. Nenhum arquivo real da raiz precisou ser substituído.

Foram usados AddressSanitizer e UndefinedBehaviorSanitizer para
detectar problemas de memória e comportamento indefinido durante
os testes. A detecção de vazamentos foi desativada porque
LeakSanitizer não funcionou sob o mecanismo de execução disponível.
Essa configuração não permite concluir que o programa esteja livre
de vazamentos.

O ambiente de execução dos testes não disponibilizou
\texttt{/dev/fuse}. Dessa forma, a avaliação não incluiu uma montagem
real nem a comunicação completa entre aplicação e kernel. Para
essa etapa posterior, o procedimento proposto é compilar o programa,
montá-lo e acessar o arquivo, por exemplo:

\begin{lstlisting}[language=bash]
mkdir -p /tmp/meufs
./passthrough-status -f /tmp/meufs
# Em outro terminal:
cat /tmp/meufs/status.txt
cat /tmp/meufs/etc/hostname
cat /tmp/meufs/status.txt
\end{lstlisting}

Esses comandos são um roteiro de demonstração, não um registro de
execução usado para produzir os resultados da próxima seção.
A desmontagem pode ser feita com \texttt{fusermount3 -u /tmp/meufs}.

\section{Resultados Parciais e Discussão}

A compilação foi concluída nas duas configurações verificadas.
Nos testes diretos, o arquivo virtual apresentou atributos de arquivo
regular com permissões 0444, aceitou leitura e recusou acesso de
escrita. A Tabela~\ref{tab:testes} apresenta os principais resultados.

\begin{table}[ht]
\centering
\caption{Resultados observados nos testes diretos dos callbacks.}
\label{tab:testes}
\begin{tabular}{p{0.31\textwidth}p{0.60\textwidth}}
\hline
Verificação & Resultado observado \\
\hline
Leitura de um byte por vez & Texto reconstruído e retorno zero no final. \\
Atualização após abrir & Abertura antiga preservada; nova abertura atualizada. \\
Criação repetida & Um incremento no primeiro sucesso; erro na repetição. \\
Listagem da raiz & Uma entrada virtual, inclusive com nome real coincidente. \\
Operações de alteração & Pedidos de escrita e alteração do caminho recusados. \\
Concorrência & 800 aberturas e leituras adicionais contabilizadas. \\
Arquivos reais & Criação, escrita, leitura e remoção funcionaram no teste. \\
\hline
\end{tabular}
\end{table}

O teste concorrente utilizou oito threads. Cada uma realizou
100 ciclos de abertura, leitura de um byte e liberação de um arquivo
real. Antes desses ciclos, havia uma abertura bem-sucedida por
\texttt{create} e dois pedidos de leitura, um deles sobre um caminho
inexistente. Ao final, foram verificados 801 registros de abertura
e 802 pedidos de leitura. Esses totais correspondem às ações
executadas pelo teste, e não a uma medição de desempenho.

A leitura do arquivo virtual foi repetida na mesma abertura depois
de alterar os contadores. O texto permaneceu igual ao obtido no
início. Após fechar e reabrir, os novos valores apareceram. Também
foram verificados tamanho zero de leitura, posição negativa e
posição além do fim do conteúdo. Isso confirma o tratamento desses
casos no callback testado.

O comportamento observado atende ao objetivo de expor informações
por meio de um arquivo gerado em memória. Entretanto, o contador
de leitura não representa todas as leituras feitas pelas aplicações:
o cache pode atender a um pedido sem chamar o programa FUSE. Também
não existe uma relação obrigatória entre um comando \texttt{cat}
e um único incremento.

Há outras limitações na versão atual. Os contadores não são
persistidos, o nome na raiz é fixo e o passthrough continua espelhando
o sistema existente. Cada contador é lido atomicamente, mas a leitura
dos três não constitui uma única transação; sob concorrência, os
valores podem corresponder a instantes próximos. Depois de gerado,
o texto de cada abertura permanece estável.

Os testes não avaliaram tempo de resposta, consumo de memória ou
comportamento sob carga prolongada. Também não verificaram todos
os detalhes herdados do exemplo original. Os resultados sustentam
as funções exercitadas diretamente, enquanto a montagem real
permanece necessária para completar a avaliação.

\section{Plano para a Entrega Final}\label{sec:plano-final}

A próxima etapa deverá transformar o relatório de contadores em um
monitor de operações. As funcionalidades descritas nesta seção são
propostas para a entrega final e ainda não fazem parte do código avaliado.

\subsection{Contagem de bytes e histórico virtual}

O arquivo \texttt{status.txt} deverá informar também o total de bytes
lidos e escritos pelos callbacks. O incremento usará o retorno positivo
de \texttt{pread} ou \texttt{pwrite}, em vez do tamanho solicitado.
Assim, leituras parciais terão seu volume efetivo contabilizado, enquanto
falhas e leituras no fim do arquivo não acrescentarão bytes. A contagem
continuará limitada às operações que chegam ao programa FUSE.

Será acrescentado o arquivo virtual \texttt{/historico.txt}, somente
para leitura, com os últimos 100 registros de criação de diretório,
abertura, leitura e escrita. Cada registro deverá informar um número
sequencial, a operação, o caminho, o resultado e a quantidade de bytes,
quando aplicável. O conteúdo dos arquivos acessados não será registrado.
A consulta poderá ser feita com \texttt{cat /tmp/meufs/historico.txt}.

Os registros serão mantidos em um vetor circular: quando atingir o
limite, o registro mais antigo será substituído. Um mutex deverá proteger
a atualização e a cópia do histórico entre threads. Cada abertura do
arquivo virtual receberá uma cópia estável, seguindo a estratégia já
usada no status. Consultas aos dois arquivos virtuais serão excluídas
do monitoramento. O histórico ficará em memória e será reiniciado
junto com o processo.

\subsection{Verificação da versão final}

A validação deverá incluir montagem real e comandos como \texttt{mkdir},
\texttt{ls}, \texttt{stat} e \texttt{cat}. Os testes diretos verificarão
a soma dos bytes retornados e a retenção dos 100 registros mais recentes
após 101 operações monitoradas. Também serão repetidos os testes de
leitura em partes, proteção contra escrita e acesso concorrente.
A entrega final deverá apresentar evidências dessa integração e comparar
os resultados com os requisitos propostos.

\section{Considerações Parciais}

A primeira etapa disponibilizou o arquivo virtual de monitoramento e
integrou consulta de atributos, listagem, abertura, leitura e liberação.
Os testes diretos confirmaram os comportamentos avaliados, incluindo as
atualizações concorrentes dos contadores. A validação pela montagem real
permanece pendente.

Essa base será reutilizada na entrega final para acrescentar a contagem
de bytes e o histórico virtual. A ampliação permitirá observar tanto os
totais acumulados quanto as operações recentes, mantendo o acesso por
comandos comuns do sistema. Os resultados apresentados nesta parcial
se referem apenas à implementação atual de \texttt{status.txt}.

\bibliographystyle{sbc}
\bibliography{referencias}
\end{document}
